
\subsection{Application Example: Palletizing LEA}
\label{subsec:PalettierenLea}

To illustrate the concepts described in Chapter~\ref{chap:Art1LEA}, this section presents an exemplary application of these concepts to a Palletizing LEA (PAL) that follows a CES-based mode of operation.

\subsubsection{Service Specification}

Table~\ref{tab:PalDienst} provides an overview of the PAL service, including its procedures, parameters, report values, and process values.

% Profile PAL service
\begin{footnotesize}
	\begin{longtable}{|
		>{\columncolor[HTML]{E0E0E0}}l
		>{\columncolor[HTML]{E0E0E0}}l
		>{\columncolor[HTML]{E0E0E0}}l
		>{\columncolor[HTML]{FFFFFF}}l |}
		\caption{Profile of an exemplary palletizing service}
		\label{tab:PalDienst}\\
		\hline

		\multicolumn{4}{|L{14.8cm}|}{\cellcolor[HTML]{E0E0E0}\textbf{Palletizing Service}}
		\\ \hline

		\multicolumn{4}{|L{14.8cm}|}{\cellcolor[HTML]{FFFFFF} This service is used to palletize bags onto pallets according to a defined packing pattern. The bags are always processed uniformly in accordance with specific order data. }
		\\ \hline

		\multicolumn{4}{|L{14.8cm}|}{\cellcolor[HTML]{E0E0E0}\textbf{Procedures}}
		\\ \hline

		\multicolumn{1}{|L{2.5cm}|}{\cellcolor[HTML]{E0E0E0}\textbf{Procedure ID}}        	&
		\multicolumn{1}{L{3.3cm}|}{\cellcolor[HTML]{E0E0E0}\textbf{Name}} 			&
		\multicolumn{2}{L{8.65cm}|}{\cellcolor[HTML]{E0E0E0}\textbf{Description}}
		\\ \hline

		\multicolumn{1}{|L{2.5cm}|}{\cellcolor[HTML]{FFFFFF} 16\#1}                     &
		\multicolumn{1}{L{3.3cm}|}{\cellcolor[HTML]{FFFFFF} CES\_Continuous (continuous)}              &
		\multicolumn{2}{L{8.65cm}|}{\cellcolor[HTML]{FFFFFF} Continuous palletizing of bags with an always identical packing pattern}
		\\ \hline

		\multicolumn{1}{|L{2.5cm}|}{\cellcolor[HTML]{FFFFFF} 16\#2}                     &
		\multicolumn{1}{L{3.3cm}|}{\cellcolor[HTML]{FFFFFF} CES\_BagQuantity (self-completing)}              &
		\multicolumn{2}{L{8.65cm}|}{\cellcolor[HTML]{FFFFFF} Palletizing a defined number of bags with an always identical packing pattern}
		\\ \hline

		\multicolumn{1}{|L{2.5cm}|}{\cellcolor[HTML]{FFFFFF} 16\#3}                     &
		\multicolumn{1}{L{3.3cm}|}{\cellcolor[HTML]{FFFFFF} CES\_PalletQuantity (self-completing)}              &
		\multicolumn{2}{L{8.65cm}|}{\cellcolor[HTML]{FFFFFF} Palletizing a defined number of pallets with an always identical packing pattern}
		\\ \hline

		\multicolumn{4}{|L{14.8cm}|}{\cellcolor[HTML]{E0E0E0}\textbf{Order-Specific Parameters (Procedure Parameters)}}
		\\ \hline

		\multicolumn{1}{|L{2.5cm}|}{\cellcolor[HTML]{E0E0E0}\textbf{Name}}        	&
		\multicolumn{1}{L{3.3cm}|}{\cellcolor[HTML]{E0E0E0}\textbf{Interface Type}} 			&
		\multicolumn{1}{L{3.17cm}|}{\cellcolor[HTML]{E0E0E0}\textbf{Procedure}} 			&
		\multicolumn{1}{L{5.3cm}|}{\cellcolor[HTML]{E0E0E0}\textbf{Description}}
		\\ \hline


		\multicolumn{1}{|L{2.5cm}|}{\cellcolor[HTML]{FFFFFF} ProductionId}        	&
		\multicolumn{1}{L{3.3cm}|}{\cellcolor[HTML]{FFFFFF} StringServParam} 			&
		\multicolumn{1}{L{3.17cm}|}{\cellcolor[HTML]{FFFFFF} 16\#1, 16\#2, 16\#3} 			&
		\multicolumn{1}{L{5.3cm}|}{\cellcolor[HTML]{FFFFFF} Order number}
		\\ \hline

		\multicolumn{1}{|L{2.5cm}|}{\cellcolor[HTML]{FFFFFF} ProductId}        	&
		\multicolumn{1}{L{3.3cm}|}{\cellcolor[HTML]{FFFFFF} DIntServParam} 			&
		\multicolumn{1}{L{3.17cm}|}{\cellcolor[HTML]{FFFFFF} 16\#1, 16\#2, 16\#3} 			&
		\multicolumn{1}{L{5.3cm}|}{\cellcolor[HTML]{FFFFFF} Identifier of the product to be palletized and key for selecting product-specific parameters}
		\\ \hline

		\multicolumn{1}{|L{2.5cm}|}{\cellcolor[HTML]{FFFFFF} BagQuantity}        	&
		\multicolumn{1}{L{3.3cm}|}{\cellcolor[HTML]{FFFFFF} DIntServParam} 			&
		\multicolumn{1}{L{3.17cm}|}{\cellcolor[HTML]{FFFFFF} 16\#2} 			&
		\multicolumn{1}{L{5.3cm}|}{\cellcolor[HTML]{FFFFFF} Number of bags to be palletized}
		\\ \hline

		\multicolumn{1}{|L{2.5cm}|}{\cellcolor[HTML]{FFFFFF} PalletQuantity}        	&
		\multicolumn{1}{L{3.3cm}|}{\cellcolor[HTML]{FFFFFF} DIntServParam} 			&
		\multicolumn{1}{L{3.17cm}|}{\cellcolor[HTML]{FFFFFF} 16\#3} 			&
		\multicolumn{1}{L{5.3cm}|}{\cellcolor[HTML]{FFFFFF} Number of pallets to be palletized}
		\\ \hline

		\multicolumn{4}{|L{14.8cm}|}{\cellcolor[HTML]{E0E0E0}\textbf{Product- and Packaging-Specific Parameters (Configuration Parameters)}}
		\\ \hline

		\multicolumn{1}{|L{2.5cm}|}{\cellcolor[HTML]{E0E0E0}\textbf{Name}}        	&
		\multicolumn{1}{L{3.3cm}|}{\cellcolor[HTML]{E0E0E0}\textbf{Interface Type}} 			&
		\multicolumn{2}{L{8.65cm}|}{\cellcolor[HTML]{E0E0E0}\textbf{Description}}
		\\ \hline

		\multicolumn{1}{|L{2.5cm}|}{\cellcolor[HTML]{FFFFFF} ProductDataSet}        	&
		\multicolumn{1}{L{3.3cm}|}{\cellcolor[HTML]{FFFFFF} ArrayServParam of PalletizerProduct-DataStruct} 			&
		\multicolumn{2}{L{8.65cm}|}{\cellcolor[HTML]{FFFFFF} Array of product-specific parameter sets; contains, e.g., the packing pattern to be used}
		\\ \hline

		\multicolumn{1}{|L{2.5cm}|}{\cellcolor[HTML]{FFFFFF} PackagingData-Set}        	&
		\multicolumn{1}{L{3.3cm}|}{\cellcolor[HTML]{FFFFFF} ArrayServParam of PalletizerPackagingDataStruct} 			&
		\multicolumn{2}{L{8.65cm}|}{\cellcolor[HTML]{FFFFFF} Array of packaging-specific parameter sets; contains, e.g., the dimensions of the pallets to be used}
		\\ \hline

		\multicolumn{4}{|L{14.8cm}|}{\cellcolor[HTML]{E0E0E0}\textbf{Design-Specific Parameters (Configuration Parameters)}}
		\\ \hline

		\multicolumn{4}{|L{14.8cm}|}{\cellcolor[HTML]{FFFFFF} None }
		\\ \hline

		\multicolumn{4}{|L{14.8cm}|}{\cellcolor[HTML]{E0E0E0}\textbf{Report Values}}
		\\ \hline

		\multicolumn{1}{|L{2.5cm}|}{\cellcolor[HTML]{E0E0E0}\textbf{Name}}        	&
		\multicolumn{1}{L{3.3cm}|}{\cellcolor[HTML]{E0E0E0}\textbf{Interface Type}} 			&
		\multicolumn{1}{L{3.17cm}|}{\cellcolor[HTML]{E0E0E0}\textbf{Procedure}} 			&
		\multicolumn{1}{L{5.3cm}|}{\cellcolor[HTML]{E0E0E0}\textbf{Description}}
		\\ \hline

		\multicolumn{1}{|L{2.5cm}|}{\cellcolor[HTML]{FFFFFF} ProductName}        	&
		\multicolumn{1}{L{3.3cm}|}{\cellcolor[HTML]{FFFFFF} StringView} 			&
		\multicolumn{1}{L{3.17cm}|}{\cellcolor[HTML]{FFFFFF} 16\#1, 16\#2, 16\#3} 			&
		\multicolumn{1}{L{5.3cm}|}{\cellcolor[HTML]{FFFFFF} Name of the currently packaged product type}
		\\ \hline

		\multicolumn{1}{|L{2.5cm}|}{\cellcolor[HTML]{FFFFFF} PalletType-Name}        	&
		\multicolumn{1}{L{3.3cm}|}{\cellcolor[HTML]{FFFFFF} StringView} 			&
		\multicolumn{1}{L{3.17cm}|}{\cellcolor[HTML]{FFFFFF} 16\#1, 16\#2, 16\#3} 			&
		\multicolumn{1}{L{5.3cm}|}{\cellcolor[HTML]{FFFFFF} Name of the currently used pallet type}
		\\ \hline

		\multicolumn{1}{|L{2.5cm}|}{\cellcolor[HTML]{FFFFFF} BagCount}        	&
		\multicolumn{1}{L{3.3cm}|}{\cellcolor[HTML]{FFFFFF} DIntView} 			&
		\multicolumn{1}{L{3.17cm}|}{\cellcolor[HTML]{FFFFFF} 16\#1, 16\#2, 16\#3} 			&
		\multicolumn{1}{L{5.3cm}|}{\cellcolor[HTML]{FFFFFF} Number of bags packaged so far}
		\\ \hline

		\multicolumn{1}{|L{2.5cm}|}{\cellcolor[HTML]{FFFFFF} PalletCount}        	&
		\multicolumn{1}{L{3.3cm}|}{\cellcolor[HTML]{FFFFFF} DIntView} 			&
		\multicolumn{1}{L{3.17cm}|}{\cellcolor[HTML]{FFFFFF} 16\#1, 16\#2, 16\#3} 			&
		\multicolumn{1}{L{5.3cm}|}{\cellcolor[HTML]{FFFFFF} Number of pallets completed so far}
		\\ \hline

		\multicolumn{4}{|L{14.8cm}|}{\cellcolor[HTML]{E0E0E0}\textbf{Process Values}}
		\\ \hline

		\multicolumn{1}{|L{2.5cm}|}{\cellcolor[HTML]{E0E0E0}\textbf{Name}}        	&
		\multicolumn{1}{L{3.3cm}|}{\cellcolor[HTML]{E0E0E0}\textbf{Interface Type}} 			&
		\multicolumn{1}{L{3.17cm}|}{\cellcolor[HTML]{E0E0E0}\textbf{Procedure}} 			&
		\multicolumn{1}{L{5.3cm}|}{\cellcolor[HTML]{E0E0E0}\textbf{Description}}
		\\ \hline

		\multicolumn{1}{|L{2.5cm}|}{\cellcolor[HTML]{FFFFFF} SuccessorClear-In}        	&
		\multicolumn{1}{L{3.3cm}|}{\cellcolor[HTML]{FFFFFF} BinProcessValueIn} 			&
		\multicolumn{1}{L{3.17cm}|}{\cellcolor[HTML]{FFFFFF} 16\#1, 16\#2, 16\#3} 			&
		\multicolumn{1}{L{5.3cm}|}{\cellcolor[HTML]{FFFFFF} Input for the clearance signal from a downstream LEA (relevant when using the LEA within a Logistics Line)}
		\\ \hline

		\multicolumn{1}{|L{2.5cm}|}{\cellcolor[HTML]{FFFFFF} Predecessor-ClearOut}        	&
		\multicolumn{1}{L{3.3cm}|}{\cellcolor[HTML]{FFFFFF} BinProcessValueIn} 			&
		\multicolumn{1}{L{3.17cm}|}{\cellcolor[HTML]{FFFFFF} 16\#1, 16\#2, 16\#3} 			&
		\multicolumn{1}{L{5.3cm}|}{\cellcolor[HTML]{FFFFFF} Output for the clearance signal to an upstream LEA (relevant when using the LEA within a Logistics Line)}
		\\ \hline


	\end{longtable}
\end{footnotesize}

\subsubsection{Mode of Operation}

The mode of operation of the PAL service follows the CES operation described in Section~\ref{sec:CES} and is illustrated in Figure~\ref{fig:ArbeitsweisePal}.

\begin{figure}[htbp]
	\centerline{\includegraphics[width=0.9\textwidth]{Inhalt/Abbildungen/99_Anhang/Spezifikation_LEA/Anwendung/Arbeitsweise_CES_PAL.png}}
	\caption{Mode of operation of an exemplary palletizing service in CES mode}
	\label{fig:ArbeitsweisePal}
\end{figure}

Initially, the service is in the IDLE state. Pre-population of the \emph{ProductDataSet} and \emph{PackagingDataSet} via the corresponding \emph{ArrayServParam} interfaces is possible. Before the service is started, the desired procedure must also be selected and all required order-specific parameters must be set, i.e., the \emph{ProductionId}, the \emph{ProductId}, the \emph{BagQuantity} (for procedure 16\#2 only), and the \emph{PalletQuantity} (for procedure 16\#3 only). Based on the configured \emph{ProductId}, the product-specific parameter set to be used is then determined from the \emph{ProductDataSet} or by querying the LOL, and selected for processing (as described in Section~\ref{subsec:Parametrierngsmech}). Using the \emph{PackagingId} contained in the selected product-specific parameter set, the packaging-specific parameter set to be used is determined and selected in the same manner (as described in Section~\ref{subsec:Parametrierngsmech}). The PAL service can then be started, and all bags belonging to the order are palletized in the EXECUTE state according to the packing pattern stored in the product parameter set. The palletizing operation is terminated when the number of bags configured via \emph{BagQuantity} is reached (applies to procedure 16\#2), when the number of pallets configured via \emph{PalletQuantity} is reached (applies to procedure 16\#3), or when the palletizing operation is ended by a \emph{Complete} command (applies to procedure 16\#1). The service can be reset to IDLE via a \emph{Reset} command.

\subsubsection{HMI Screen}

Figure~\ref{fig:HmiPal} shows the HMI screen of the PAL. It enables the entry of order-specific parameters, control of the LEA service, and the display of report values and process values via dynamic HMI objects in accordance with \cite{PNO.2025.Part2}. In addition, it contains an image of the PAL as a static HMI object. This image is stored in the MTP attachment of the PAL as shown in Figure~\ref{fig:HmiPalMtpAnhang}. It is assigned the ECLASS reference \enquote{36521201}, which corresponds to the semantic reference for a palletizer. A LOL integrating this LEA can therefore decide whether to use the image from the MTP attachment or an image of a palletizer that may be available in its graphics library.

\begin{figure}[htbp]
	\centerline{\includegraphics[width=1.0\textwidth]{Inhalt/Abbildungen/99_Anhang/Spezifikation_LEA/Anwendung/HMI_PAL_PVs.png}}
	\caption{HMI screen of a palletizing LEA}
	\label{fig:HmiPal}
\end{figure}

\begin{figure}[htbp]
	\centerline{\includegraphics[width=0.5\textwidth]{Inhalt/Abbildungen/99_Anhang/Spezifikation_LEA/Anwendung/MTP_HMI_PAL.drawio.png}}
	\caption{Image file of a palletizing LEA within the MTP}
	\label{fig:HmiPalMtpAnhang}
\end{figure}

\subsection{Application Example: Stretch Hood LEA}
\label{subsec:StretchenLea}

To illustrate the concepts described in Chapter~\ref{chap:Art1LEA}, this section presents an exemplary application of these concepts to a Stretch Hood LEA (SH), which can operate in either CES-based or SES-based mode.

\subsubsection{Service Specification}

Table~\ref{tab:ShDienst} provides an overview of the SH service, including its procedures, parameters, report values, and process values.

% Profile SH service
\begin{footnotesize}
	\begin{longtable}{|
		>{\columncolor[HTML]{E0E0E0}}l
		>{\columncolor[HTML]{E0E0E0}}l
		>{\columncolor[HTML]{E0E0E0}}l
		>{\columncolor[HTML]{FFFFFF}}l |}
		\caption{Profile of an exemplary stretch hood service}
		\label{tab:ShDienst}\\
		\hline

		\multicolumn{4}{|L{14.8cm}|}{\cellcolor[HTML]{E0E0E0}\textbf{Stretch Hood Service}}
		\\ \hline

		\multicolumn{4}{|L{14.8cm}|}{\cellcolor[HTML]{FFFFFF} Depending on its configuration, this service can apply film hoods to pallets containing bags or octabins, or insert film liners into empty octabins. Depending on the plant context, the SH can operate in an order-oriented manner or fulfill its various tasks on demand, potentially for different packaging processes.}
		\\ \hline

		\multicolumn{4}{|L{14.8cm}|}{\cellcolor[HTML]{E0E0E0}\textbf{Procedures}}
		\\ \hline

		\multicolumn{1}{|L{2.5cm}|}{\cellcolor[HTML]{E0E0E0}\textbf{Procedure ID}}        	&
		\multicolumn{1}{L{3.3cm}|}{\cellcolor[HTML]{E0E0E0}\textbf{Name}} 			&
		\multicolumn{2}{L{8.65cm}|}{\cellcolor[HTML]{E0E0E0}\textbf{Description}}
		\\ \hline

		\multicolumn{1}{|L{2.5cm}|}{\cellcolor[HTML]{FFFFFF} 16\#1}                     &
		\multicolumn{1}{L{3.3cm}|}{\cellcolor[HTML]{FFFFFF} CES\_Continuous (continuous)}              &
		\multicolumn{2}{L{8.65cm}|}{\cellcolor[HTML]{FFFFFF} Continuous execution of \underline{one} of the above tasks in an always identical manner according to order data}
		\\ \hline

		\multicolumn{1}{|L{2.5cm}|}{\cellcolor[HTML]{FFFFFF} 16\#2}                     &
		\multicolumn{1}{L{3.3cm}|}{\cellcolor[HTML]{FFFFFF} CES\_PalletQuantity (self-completing)}              &
		\multicolumn{2}{L{8.65cm}|}{\cellcolor[HTML]{FFFFFF} Execution of \underline{one} of the above tasks in an always identical manner on a defined number of LOs according to order data}
		\\ \hline

		\multicolumn{1}{|L{2.5cm}|}{\cellcolor[HTML]{FFFFFF} 16\#3}                     &
		\multicolumn{1}{L{3.3cm}|}{\cellcolor[HTML]{FFFFFF} SES\_Continuous (continuous)}              &
		\multicolumn{2}{L{8.65cm}|}{\cellcolor[HTML]{FFFFFF} Flexible execution of \underline{different} tasks from above depending on the arriving LO}
		\\ \hline

		\multicolumn{4}{|L{14.8cm}|}{\cellcolor[HTML]{E0E0E0}\textbf{Order-Specific Parameters (Procedure Parameters)}}
		\\ \hline

		\multicolumn{1}{|L{2.5cm}|}{\cellcolor[HTML]{E0E0E0}\textbf{Name}}        	&
		\multicolumn{1}{L{3.3cm}|}{\cellcolor[HTML]{E0E0E0}\textbf{Interface Type}} 			&
		\multicolumn{1}{L{3.17cm}|}{\cellcolor[HTML]{E0E0E0}\textbf{Procedure}} 			&
		\multicolumn{1}{L{5.3cm}|}{\cellcolor[HTML]{E0E0E0}\textbf{Description}}
		\\ \hline


		\multicolumn{1}{|L{2.5cm}|}{\cellcolor[HTML]{FFFFFF} ProductionId}        	&
		\multicolumn{1}{L{3.3cm}|}{\cellcolor[HTML]{FFFFFF} StringServParam} 			&
		\multicolumn{1}{L{3.17cm}|}{\cellcolor[HTML]{FFFFFF} 16\#1, 16\#2} 			&
		\multicolumn{1}{L{5.3cm}|}{\cellcolor[HTML]{FFFFFF} Order number}
		\\ \hline

		\multicolumn{1}{|L{2.5cm}|}{\cellcolor[HTML]{FFFFFF} ProductId}        	&
		\multicolumn{1}{L{3.3cm}|}{\cellcolor[HTML]{FFFFFF} DIntServParam} 			&
		\multicolumn{1}{L{3.17cm}|}{\cellcolor[HTML]{FFFFFF} 16\#1, 16\#2} 			&
		\multicolumn{1}{L{5.3cm}|}{\cellcolor[HTML]{FFFFFF} Identifier of the product to be processed and key for selecting product-specific parameters}
		\\ \hline

		\multicolumn{1}{|L{2.5cm}|}{\cellcolor[HTML]{FFFFFF} PalletQuantity}        	&
		\multicolumn{1}{L{3.3cm}|}{\cellcolor[HTML]{FFFFFF} DIntServParam} 			&
		\multicolumn{1}{L{3.17cm}|}{\cellcolor[HTML]{FFFFFF} 16\#2} 			&
		\multicolumn{1}{L{5.3cm}|}{\cellcolor[HTML]{FFFFFF} Number of LOs to be processed}
		\\ \hline

		\multicolumn{4}{|L{14.8cm}|}{\cellcolor[HTML]{E0E0E0}\textbf{Product- and Packaging-Specific Parameters (Configuration Parameters)}}
		\\ \hline

		\multicolumn{1}{|L{2.5cm}|}{\cellcolor[HTML]{E0E0E0}\textbf{Name}}        	&
		\multicolumn{1}{L{3.3cm}|}{\cellcolor[HTML]{E0E0E0}\textbf{Interface Type}} 			&
		\multicolumn{2}{L{8.65cm}|}{\cellcolor[HTML]{E0E0E0}\textbf{Description}}
		\\ \hline

		\multicolumn{1}{|L{2.5cm}|}{\cellcolor[HTML]{FFFFFF} ProductDataSet}        	&
		\multicolumn{1}{L{3.3cm}|}{\cellcolor[HTML]{FFFFFF} ArrayServParam of StretchHoodProductDataStruct} 			&
		\multicolumn{2}{L{8.65cm}|}{\cellcolor[HTML]{FFFFFF} Array of product-specific parameter sets; contains, e.g., the stretch parameters to be used}
		\\ \hline

		\multicolumn{1}{|L{2.5cm}|}{\cellcolor[HTML]{FFFFFF} PackagingData-Set}        	&
		\multicolumn{1}{L{3.3cm}|}{\cellcolor[HTML]{FFFFFF} ArrayServParam of StretchHoodPackagingDataStruct} 			&
		\multicolumn{2}{L{8.65cm}|}{\cellcolor[HTML]{FFFFFF} Array of packaging-specific parameter sets; contains, e.g., the properties of the stretch film to be used}
		\\ \hline

		\multicolumn{4}{|L{14.8cm}|}{\cellcolor[HTML]{E0E0E0}\textbf{Design-Specific Parameters (Configuration Parameters)}}
		\\ \hline

		\multicolumn{1}{|L{2.5cm}|}{\cellcolor[HTML]{E0E0E0}\textbf{Name}}        	&
		\multicolumn{1}{L{3.3cm}|}{\cellcolor[HTML]{E0E0E0}\textbf{Interface Type}} 			&
		\multicolumn{2}{L{8.65cm}|}{\cellcolor[HTML]{E0E0E0}\textbf{Description}}
		\\ \hline

		\multicolumn{1}{|L{2.5cm}|}{\cellcolor[HTML]{FFFFFF} FoilTypeId}        	&
		\multicolumn{1}{L{3.3cm}|}{\cellcolor[HTML]{FFFFFF} DIntServParam} 			&
		\multicolumn{2}{L{8.65cm}|}{\cellcolor[HTML]{FFFFFF} Identifier of the loaded film type}
		\\ \hline

		\multicolumn{4}{|L{14.8cm}|}{\cellcolor[HTML]{E0E0E0}\textbf{Report Values}}
		\\ \hline

		\multicolumn{1}{|L{2.5cm}|}{\cellcolor[HTML]{E0E0E0}\textbf{Name}}        	&
		\multicolumn{1}{L{3.3cm}|}{\cellcolor[HTML]{E0E0E0}\textbf{Interface Type}} 			&
		\multicolumn{1}{L{3.17cm}|}{\cellcolor[HTML]{E0E0E0}\textbf{Procedure}} 			&
		\multicolumn{1}{L{5.3cm}|}{\cellcolor[HTML]{E0E0E0}\textbf{Description}}
		\\ \hline

		\multicolumn{1}{|L{2.5cm}|}{\cellcolor[HTML]{FFFFFF} ProductName}        	&
		\multicolumn{1}{L{3.3cm}|}{\cellcolor[HTML]{FFFFFF} StringView} 			&
		\multicolumn{1}{L{3.17cm}|}{\cellcolor[HTML]{FFFFFF} 16\#1, 16\#2, 16\#3} 			&
		\multicolumn{1}{L{5.3cm}|}{\cellcolor[HTML]{FFFFFF} Name of the currently packaged product type}
		\\ \hline

		\multicolumn{1}{|L{2.5cm}|}{\cellcolor[HTML]{FFFFFF} FoilTypeName}        	&
		\multicolumn{1}{L{3.3cm}|}{\cellcolor[HTML]{FFFFFF} StringView} 			&
		\multicolumn{1}{L{3.17cm}|}{\cellcolor[HTML]{FFFFFF} 16\#1, 16\#2, 16\#3} 			&
		\multicolumn{1}{L{5.3cm}|}{\cellcolor[HTML]{FFFFFF} Name of the currently used film type}
		\\ \hline

		\multicolumn{1}{|L{2.5cm}|}{\cellcolor[HTML]{FFFFFF} PalletCount}        	&
		\multicolumn{1}{L{3.3cm}|}{\cellcolor[HTML]{FFFFFF} DIntView} 			&
		\multicolumn{1}{L{3.17cm}|}{\cellcolor[HTML]{FFFFFF} 16\#1, 16\#2, 16\#3} 			&
		\multicolumn{1}{L{5.3cm}|}{\cellcolor[HTML]{FFFFFF} Number of LOs processed so far}
		\\ \hline

		\multicolumn{4}{|L{14.8cm}|}{\cellcolor[HTML]{E0E0E0}\textbf{Process Values}}
		\\ \hline

		\multicolumn{1}{|L{2.5cm}|}{\cellcolor[HTML]{E0E0E0}\textbf{Name}}        	&
		\multicolumn{1}{L{3.3cm}|}{\cellcolor[HTML]{E0E0E0}\textbf{Interface Type}} 			&
		\multicolumn{1}{L{3.17cm}|}{\cellcolor[HTML]{E0E0E0}\textbf{Procedure}} 			&
		\multicolumn{1}{L{5.3cm}|}{\cellcolor[HTML]{E0E0E0}\textbf{Description}}
		\\ \hline

		\multicolumn{1}{|L{2.5cm}|}{\cellcolor[HTML]{FFFFFF} SuccessorClear-In}        	&
		\multicolumn{1}{L{3.3cm}|}{\cellcolor[HTML]{FFFFFF} BinProcessValueIn} 			&
		\multicolumn{1}{L{3.17cm}|}{\cellcolor[HTML]{FFFFFF} 16\#1, 16\#2, 16\#3} 			&
		\multicolumn{1}{L{5.3cm}|}{\cellcolor[HTML]{FFFFFF} Input for the clearance signal from a downstream LEA (relevant when using the LEA within a Logistics Line)}
		\\ \hline

		\multicolumn{1}{|L{2.5cm}|}{\cellcolor[HTML]{FFFFFF} Predecessor-ClearOut}        	&
		\multicolumn{1}{L{3.3cm}|}{\cellcolor[HTML]{FFFFFF} BinProcessValueIn} 			&
		\multicolumn{1}{L{3.17cm}|}{\cellcolor[HTML]{FFFFFF} 16\#1, 16\#2, 16\#3} 			&
		\multicolumn{1}{L{5.3cm}|}{\cellcolor[HTML]{FFFFFF} Output for the clearance signal to an upstream LEA (relevant when using the LEA within a Logistics Line)}
		\\ \hline


	\end{longtable}
\end{footnotesize}

\subsubsection{Mode of Operation}

The mode of operation of the two CES procedures is similar to that of the PAL service. Therefore, this section focuses on the SES procedure of the SH service (16\#3). It follows the SES operation described in Section~\ref{sec:SES} and is illustrated in Figure~\ref{fig:ArbeitsweiseSh}.

\begin{figure}[htbp]
	\centerline{\includegraphics[width=0.85\textwidth]{Inhalt/Abbildungen/99_Anhang/Spezifikation_LEA/Anwendung/Arbeitsweise_SES_SH.png}}
	\caption{Mode of operation of an exemplary stretch hood service in SES mode}
	\label{fig:ArbeitsweiseSh}
\end{figure}

Initially, the service is in the IDLE state. Pre-population of the \emph{ProductDataSet} and \emph{PackagingDataSet} via the corresponding \emph{ArrayServParam} interfaces is possible. In addition, the LEA must be loaded with the desired film type, and its ID must be set at the design-specific parameter \emph{FoilTypeId}. Using this ID, the associated packaging-specific parameters can be selected from the \emph{PackagingDataSet} or requested from the LO's parameter management (as described in Section~\ref{subsec:Parametrierngsmech}).

Before the service is started, the desired procedure must also be selected (here: 16\#3). The procedure is then started and transitions to the PAUSED state. When an LO arrives, a state transition to RESUMING takes place, and the \emph{ProductId} as well as the \emph{LogisticsObjectStatus} of the LO are determined (see also Chapter~\ref{chap:Art3LA}, \emph{ProductId} and \emph{LogisticsObjectStatus} as part of the interface of a transport service). Based on this, the product-specific parameter set to be used is then determined from the \emph{ProductDataSet} or by querying the LOL, and selected for processing (as described in Section~\ref{subsec:Parametrierngsmech}). A state transition to EXECUTE follows. As the basis for processing, the \emph{PackagingId} contained in the product data set is compared with the \emph{FoilTypeId} configured at the LEA. If they match, the arrived LO is processed according to the configured parameters. If they do not match, the pallet cannot be processed by the LEA and is rejected (see also Section~\ref{subsec:Prozessmodell}, TransportDecline).

The service then transitions back to the PAUSED state and waits for the next incoming LO. If no further LOs are to be processed, the service can be terminated via a \emph{Complete} command and reset to IDLE via a \emph{Reset} command.

\subsubsection{HMI Screen}

Figure~\ref{fig:HmiSh} shows the HMI screen of the SH. It enables the entry of order-specific parameters, control of the LEA service, and the display of report values and process values via dynamic HMI objects in accordance with \cite{PNO.2025.Part2}. In addition, it contains an image of the SH as a static HMI object. This image is stored in the MTP attachment of the SH as shown in Figure~\ref{fig:HmiShMtpAnhang}. Since no suitable ECLASS reference exists for stretch hood machines, the reference \enquote{90900001} is used. A LOL integrating this LEA must therefore use the image provided in the MTP attachment.

\begin{figure}[htbp]
	\centerline{\includegraphics[width=1.0\textwidth]{Inhalt/Abbildungen/99_Anhang/Spezifikation_LEA/Anwendung/HMI_SH_PVs.png}}
	\caption{HMI screen of a stretch hooding and film lining LEA}
	\label{fig:HmiSh}
\end{figure}

\begin{figure}[htbp]
	\centerline{\includegraphics[width=0.5\textwidth]{Inhalt/Abbildungen/99_Anhang/Spezifikation_LEA/Anwendung/MTP_HMI_SH.drawio.png}}
	\caption{Image file of a stretch hood LEA within the MTP}
	\label{fig:HmiShMtpAnhang}
\end{figure}

\subsection{Application Example: Bag Filling Line}
\label{subsec:LLAnwendung}

To illustrate the concepts described in Chapter~\ref{chap:Art2LL}, this section presents an exemplary application of the choreography concept to the bag filling line of the use case described in Section~\ref{subsec:Anwendungsbeispiel}.

\subsubsection{Introduction to the Use Case}

The bag filling line consists of three LEAs — a Form-Fill-Seal machine (FFS), a conveyor (CONV), and a palletizer (PAL) — each automated with a service in CES mode in accordance with Chapter~\ref{chap:Art1LEA}. For simplicity, it is assumed that these services feature exclusively the procedures, parameters, process values, and report values shown in Table~\ref{tab:BspDiensteLogistiklinie}.

% Services of the application example
\begin{footnotesize}
	\begin{longtable}{|
		>{\columncolor[HTML]{E0E0E0}}l
		>{\columncolor[HTML]{E0E0E0}}l
		>{\columncolor[HTML]{E0E0E0}}l
		>{\columncolor[HTML]{E0E0E0}}l
		>{\columncolor[HTML]{E0E0E0}}l |}
		\caption{Services of the exemplary Logistics Line including their procedures, parameters, and process values}
		\label{tab:BspDiensteLogistiklinie}\\
		\hline

		\multicolumn{1}{|L{2cm}|}{\cellcolor[HTML]{E0E0E0}\textbf{Service (LEA)}}        &
		\multicolumn{1}{L{3.5cm}|}{\cellcolor[HTML]{E0E0E0}\textbf{Procedure (ID)}}        &
		\multicolumn{1}{L{2.3cm}|}{\cellcolor[HTML]{E0E0E0} \textbf{Parameters}} &
		\multicolumn{1}{L{3.5cm}|}{\cellcolor[HTML]{E0E0E0} \textbf{Process Value}} &
		\multicolumn{1}{L{2.3cm}|}{\cellcolor[HTML]{E0E0E0} \textbf{Report Values}}
		\\ \hline

        \multicolumn{1}{|L{2cm}|}{\cellcolor[HTML]{FFFFFF} BagFilling }        &
		\multicolumn{1}{L{3.5cm}|}{\cellcolor[HTML]{FFFFFF} Continuous (16\#1)}        &
		\multicolumn{1}{L{2.3cm}|}{\cellcolor[HTML]{FFFFFF}  	ProductId  } &
		\multicolumn{1}{L{3.5cm}|}{\cellcolor[HTML]{FFFFFF} SuccessorClearIn}  &
		\multicolumn{1}{L{2.3cm}|}{\cellcolor[HTML]{FFFFFF} BagsCount}
		\\ \cline{2-5}

		\multicolumn{1}{|L{2cm}|}{\cellcolor[HTML]{FFFFFF} (FFS)}     &
		\multicolumn{1}{L{3.5cm}|}{\cellcolor[HTML]{FFFFFF} BagQuantity (16\#2)}        &
		\multicolumn{1}{L{2.3cm}|}{\cellcolor[HTML]{FFFFFF}  	ProductId  } &
		\multicolumn{1}{L{3.5cm}|}{\cellcolor[HTML]{FFFFFF} SuccessorClearIn} &
		\multicolumn{1}{L{2.3cm}|}{\cellcolor[HTML]{FFFFFF} BagsCount}
		\\ \cline{3-3}

		\multicolumn{1}{|L{2cm}|}{\cellcolor[HTML]{FFFFFF} }     &
		\multicolumn{1}{L{3.5cm}|}{\cellcolor[HTML]{FFFFFF} }        &
		\multicolumn{1}{L{2.3cm}|}{\cellcolor[HTML]{FFFFFF}  	BagQuantity  } &
		\multicolumn{1}{L{3.5cm}|}{\cellcolor[HTML]{FFFFFF} }    &
		\multicolumn{1}{L{2.3cm}|}{\cellcolor[HTML]{FFFFFF} }
		\\ \hline

		\multicolumn{1}{|L{2cm}|}{\cellcolor[HTML]{FFFFFF} Conveying }        &
		\multicolumn{1}{L{3.5cm}|}{\cellcolor[HTML]{FFFFFF} Continuous (16\#1)}        &
		\multicolumn{1}{L{2.3cm}|}{\cellcolor[HTML]{FFFFFF}  	Speed  } &
		\multicolumn{1}{L{3.5cm}|}{\cellcolor[HTML]{FFFFFF} PredecessorClearOut} &
		\multicolumn{1}{L{2.3cm}|}{\cellcolor[HTML]{FFFFFF} - }
		\\ \cline{4-4}

		\multicolumn{1}{|L{2cm}|}{\cellcolor[HTML]{FFFFFF} (CONV)}        &
		\multicolumn{1}{L{3.5cm}|}{\cellcolor[HTML]{FFFFFF} }        &
		\multicolumn{1}{L{2.3cm}|}{\cellcolor[HTML]{FFFFFF}  	 } &
		\multicolumn{1}{L{3.5cm}|}{\cellcolor[HTML]{FFFFFF} SuccessorClearIn} &
		\multicolumn{1}{L{2.3cm}|}{\cellcolor[HTML]{FFFFFF} }
		\\ \hline

		\multicolumn{1}{|L{2cm}|}{\cellcolor[HTML]{FFFFFF} Palletizing }        &
		\multicolumn{1}{L{3.5cm}|}{\cellcolor[HTML]{FFFFFF} Continuous (16\#1)}        &
		\multicolumn{1}{L{2.3cm}|}{\cellcolor[HTML]{FFFFFF}  	ProductId  } &
		\multicolumn{1}{L{3.5cm}|}{\cellcolor[HTML]{FFFFFF} PredecessorClearOut} &
		\multicolumn{1}{L{2.3cm}|}{\cellcolor[HTML]{FFFFFF} BagsCount}
		\\ \cline{5-5}

		\multicolumn{1}{|L{2cm}|}{\cellcolor[HTML]{FFFFFF} (PAL)}        &
		\multicolumn{1}{L{3.5cm}|}{\cellcolor[HTML]{FFFFFF} }        &
		\multicolumn{1}{L{2.3cm}|}{\cellcolor[HTML]{FFFFFF}  	  } &
		\multicolumn{1}{L{3.5cm}|}{\cellcolor[HTML]{FFFFFF} } &
		\multicolumn{1}{L{2.3cm}|}{\cellcolor[HTML]{FFFFFF} PalletsCount}
		\\ \cline{2-5}

		\multicolumn{1}{|L{2cm}|}{\cellcolor[HTML]{FFFFFF} }     &
		\multicolumn{1}{L{3.5cm}|}{\cellcolor[HTML]{FFFFFF} BagQuantity (16\#2)}        &
		\multicolumn{1}{L{2.3cm}|}{\cellcolor[HTML]{FFFFFF}  	ProductId  } &
		\multicolumn{1}{L{3.5cm}|}{\cellcolor[HTML]{FFFFFF} PredecessorClearOut} &
		\multicolumn{1}{L{2.3cm}|}{\cellcolor[HTML]{FFFFFF} BagsCount}
		\\ \cline{3-3} \cline{5-5}

		\multicolumn{1}{|L{2cm}|}{\cellcolor[HTML]{FFFFFF} }     &
		\multicolumn{1}{L{3.5cm}|}{\cellcolor[HTML]{FFFFFF} }        &
		\multicolumn{1}{L{2.3cm}|}{\cellcolor[HTML]{FFFFFF}  	BagQuantity  } &
		\multicolumn{1}{L{3.5cm}|}{\cellcolor[HTML]{FFFFFF} } &
		\multicolumn{1}{L{2.3cm}|}{\cellcolor[HTML]{FFFFFF} PalletsCount}
		\\ \cline{2-5}

		\multicolumn{1}{|L{2cm}|}{\cellcolor[HTML]{FFFFFF} }     &
		\multicolumn{1}{L{3.5cm}|}{\cellcolor[HTML]{FFFFFF} PalletQuantity (16\#3)}        &
		\multicolumn{1}{L{2.3cm}|}{\cellcolor[HTML]{FFFFFF}  	ProductId  } &
		\multicolumn{1}{L{3.5cm}|}{\cellcolor[HTML]{FFFFFF} PredecessorClearOut} &
		\multicolumn{1}{L{2.3cm}|}{\cellcolor[HTML]{FFFFFF} BagsCount}
		\\ \cline{3-3} \cline{5-5}

		\multicolumn{1}{|L{2cm}|}{\cellcolor[HTML]{FFFFFF} }     &
		\multicolumn{1}{L{3.5cm}|}{\cellcolor[HTML]{FFFFFF} }        &
		\multicolumn{1}{L{2.3cm}|}{\cellcolor[HTML]{FFFFFF}  	PalletQuantity  } &
		\multicolumn{1}{L{3.5cm}|}{\cellcolor[HTML]{FFFFFF} } &
		\multicolumn{1}{L{2.3cm}|}{\cellcolor[HTML]{FFFFFF} PalletsCount}
		\\ \hline

	\end{longtable}
\end{footnotesize}

The choreography configuration in terms of Configurable Logic and Configurable Communication depends substantially on the desired behavior of the Logistics Line. Table~\ref{tab:SzenarienLL} shows exemplary application scenarios that are used as the basis here. In other application contexts, different scenarios and thus a different choreography configuration may be required.

% Application scenarios of the bag filling line
\begin{footnotesize}
	\begin{longtable}{|
		>{\columncolor[HTML]{E0E0E0}}l
		>{\columncolor[HTML]{FFFFFF}}l |}
		\caption{Application scenarios of the exemplary bag filling line}
		\label{tab:SzenarienLL}\\
		\hline

		\multicolumn{1}{|L{2cm}|}{\cellcolor[HTML]{E0E0E0}\textbf{Scenario}}        &
		\multicolumn{1}{L{12.63cm}|}{\cellcolor[HTML]{E0E0E0} \textbf{Description}}
		\\ \hline

        \multicolumn{1}{|L{2cm}|}{\cellcolor[HTML]{FFFFFF} Scenario 1 -\newline Start-Up}        &
		\multicolumn{1}{L{12.63cm}|}{\cellcolor[HTML]{FFFFFF} The line must be started up in an orderly manner from back to front, i.e., from the PAL to the FFS. The service of an LEA must therefore be started as soon as the service of the downstream LEA is running.}
		\\ \hline

		\multicolumn{1}{|L{2cm}|}{\cellcolor[HTML]{FFFFFF} Scenario 2 -\newline Drain}        &
		\multicolumn{1}{L{12.63cm}|}{\cellcolor[HTML]{FFFFFF} The line must be drained in an orderly manner from front to back, i.e., from the FFS to the PAL. The service of an LEA must therefore be stopped as soon as the service of the upstream LEA has been stopped.}
		\\ \hline

		\multicolumn{1}{|L{2cm}|}{\cellcolor[HTML]{FFFFFF} Scenario 3 -\newline Exception Handling}        &
		\multicolumn{1}{L{12.63cm}|}{\cellcolor[HTML]{FFFFFF} When an error occurs in the packaging line, all LEAs must transition to a fault state. The service of an LEA must therefore switch to a fault state as soon as another LEA service or an existing LOL indicates a fault state.}
		\\ \hline

		\multicolumn{1}{|L{2cm}|}{\cellcolor[HTML]{FFFFFF} Scenario 4 -\newline Unhold}        &
		\multicolumn{1}{L{12.63cm}|}{\cellcolor[HTML]{FFFFFF} If the error is non-critical (service state: HELD), it must be possible to restart the Logistics Line after the fault has been resolved. The service of an LEA must therefore be restarted as soon as the service of the downstream LEA is running again.}
		\\ \hline

		\multicolumn{1}{|L{2cm}|}{\cellcolor[HTML]{FFFFFF} Scenario 5 -\newline Reset}        &
		\multicolumn{1}{L{12.63cm}|}{\cellcolor[HTML]{FFFFFF} If a critical error has occurred or the filling process has been completed, it must be possible to reset the Logistics Line to its initial state IDLE. The service of an LEA must therefore be reset as soon as the service of the downstream LEA has been reset.}
		\\ \hline



	\end{longtable}
\end{footnotesize}

\subsubsection{Choreography Configuration — Vertical Integration}

\textbf{Lead Service:} As described in Section~\ref{subsec:LineSchnittstelle}, a Lead Service provides a service interface for the functionality of a Logistics Line, which can be implemented as an Internal Lead or External Lead. In this case, an External Lead is chosen, since none of the three participating services can meaningfully represent the line functionality. Considering Scenario 1 (Start-Up), for example, it becomes clear that the impulse to start up must be given at the PAL service, while the confirmation of a successful start-up is available at the FFS service. An Internal Lead is therefore not suitable. The External Lead Service can run, for instance, alongside the \emph{Palletizing} service on the PAL's controller, since these typically serve as the master controllers for logistics lines today, but it may also run on a separate controller. The Lead Service is exemplarily equipped with five generic procedures that, in accordance with the behavioral rules of the Configurable Logic, lead to various procedure settings of the Follower Services.

\textbf{Parameters:} In addition, the parameters shown in Table~\ref{tab:BspDiensteLogistiklinie} must be configured. It is advisable to configure parameters at as few LEAs as possible in order to simplify the line interface. In the present case, all parameters of the \emph{Palletizing} service and the \emph{Speed} parameter of the \emph{Conveying} service are provided at the line interface. The parameters \emph{ProductId} and \emph{BagQuantity} are passed internally within the choreography from the PAL to the FFS LEA.

\textbf{Report Values:} The report values of the Logistics Line are the \emph{BagsCount} and \emph{PalletsCount} report values of the PAL, which provide the number of completed bags and pallets of the Logistics Line respectively.

\textbf{Process Values:} The available process values are in this case processed exclusively within the choreography and are not exposed at the line interface.

\textbf{HMI Screens:} The line HMI screen is composed of three LEA images as CustomSymbols, which are obtained from the LEA MTPs using the ContextReference mechanism in combination with the CustomSymbols mechanism. In addition, dynamic \emph{VisualObjects} are provided for the \emph{ServiceControl} interfaces of the Lead Service and the three LEA services, as well as for the parameters and report values of the line interface. These reference the associated \emph{DataAssemblies} in the attached LEA MTPs via \emph{RC HasExternalMtpContext}. Figure~\ref{fig:LineBedienbild} shows the line HMI screen. In addition, the three HMI screens of the individual LEAs are incorporated into the HMI screen hierarchy of the Logistics Line via \emph{ReferencedPicture}.

\begin{figure}[htbp]
	\centerline{\includegraphics[width=1.0\textwidth]{Inhalt/Abbildungen/07_Linen-Automatisierung/Line_STARTING_RV2.png}}
	\caption{Exemplary line HMI screen of a bag filling line}
	\label{fig:LineBedienbild}
\end{figure}

The elements of the line interface and line HMI screen described above are represented in a Composed MTP in accordance with the mechanisms described in Section~\ref{sec:KonzeptComposedMTP}.

\subsubsection{Choreography Configuration — Horizontal Integration}

Procedural, parametrizing, and interlocking relations, as well as constants, are defined between the three LEAs of the Logistics Line.

\textbf{Procedural Relations:} Figure~\ref{fig:ProzRelationen} shows an exemplary excerpt from a matrix describing the procedural relations between the LEAs of the bag filling line, specifically the relations between FFS and CONV.

\begin{figure}[htbp]
	\centerline{\includegraphics[width=0.7\textwidth]{Inhalt/Abbildungen/07_Linen-Automatisierung/Proz_Relationen.png}}
	\caption{Exemplary procedural relations of a bag filling line}
	\label{fig:ProzRelationen}
\end{figure}

The rows represent the system variables that the \emph{Conveying} service can make available to other services. The columns represent the choreography interface of the \emph{BagFilling} service, which can be actuated based on the variables of other LEAs. The markings within the matrix represent the behavioral rules of the Configurable Logic. For example, the rule highlighted in green indicates that the EXECUTE state of the \emph{Conveying} service triggers the \emph{Start} command at the \emph{BagFilling} service. This rule contributes to the orderly start-up of the line (Scenario 1). In such a direct relation, a system variable of one service has an immediate effect on the interface of another service. In the case of Scenario 3 \enquote{Fault} (Figure~\ref{fig:ProzRelationen}, red markings), it additionally occurs that two system variables act upon the same variable of the choreography interface. In this case, pre-processing of the system variables is required (Figure~\ref{fig:ProzRelationen}, \enquote{Pre-processing}), the result of which is written to the choreography interface. In this case, the pre-processing is a logical OR (\enquote{O}). All other arithmetic and logical functions can likewise be used for such pre-processing. As demonstrated here using the \emph{Conveying} and \emph{BagFilling} services as an example, all other behavioral rules between the three services of the bag filling line and the External Lead can be realized in the same way.

Furthermore, it must be specified how the procedure setting of the External Lead is to be propagated to the LEA services. In this example, a self-completing procedure 1 for packaging a defined number of pallets is configured at the Lead Service. This results in the \emph{BagQuantity} procedure being set at the \emph{BagFilling} service, the \emph{Continuous} procedure at the \emph{Conveying} service, and the \emph{PalletQuantity} procedure at the \emph{Palletizing} service.

\textbf{Parametrizing Relations:} Figure~\ref{fig:ParamRelationen} shows an exemplary excerpt from a matrix describing the parametrizing relations between the LEAs of the bag filling line, specifically the relations between PAL and FFS.

\begin{figure}[htbp]
	\centerline{\includegraphics[width=0.4\textwidth]{Inhalt/Abbildungen/07_Linen-Automatisierung/Param_Relationen.png}}
	\caption{Exemplary parametrizing relations of a bag filling line}
	\label{fig:ParamRelationen}
\end{figure}

As described previously in the section \enquote{Choreography Configuration — Vertical Integration}, the \emph{BagFilling} service obtains the parameters it requires (i.e., \emph{ProductId} and \emph{BagQuantity}) from the PAL. In the present case, no pre-processing is required for this.

\textbf{Interlocking Relations:} Figure~\ref{fig:InterlockRelationen} shows an exemplary excerpt from a matrix describing the interlocking relations between the LEAs of the bag filling line, specifically the relations between PAL and CONV.

\begin{figure}[htbp]
	\centerline{\includegraphics[width=0.5\textwidth]{Inhalt/Abbildungen/07_Linen-Automatisierung/Interlock_Relationen.png}}
	\caption{Exemplary interlocking relations of a bag filling line}
	\label{fig:InterlockRelationen}
\end{figure}

In the present use case, the available process values for transmitting the clearance signals are forwarded as interlocking relations without pre-processing.

\textbf{Constants:} Figure~\ref{fig:ConstRelationen} shows an exemplary excerpt from a matrix describing the definition of constants for the bag filling line.

\begin{figure}[htbp]
	\centerline{\includegraphics[width=0.85\textwidth]{Inhalt/Abbildungen/07_Linen-Automatisierung/Const_Relationen.png}}
	\caption{Exemplary constants of a bag filling line}
	\label{fig:ConstRelationen}
\end{figure}

The ability to define constants is used in the case of the bag filling line to define the operating mode of the LEA services in accordance with Section~\ref{subsec:LineSchnittstelle}. For the three LEA services \emph{BagFilling}, \emph{Conveying}, and \emph{Palletizing}, the \emph{Access Mode} variable is set to a constant value of \enquote{1}. The LEA controller interprets this as the \emph{Automatic Internal} operating mode. For the Lead Service, the \emph{Access Mode} is set to a constant value of \enquote{2}. The LEA controller interprets this as the \emph{Operator} operating mode. Furthermore, constants can be used, for example, to set parameters such as the \emph{Speed} variable of the CONV to a constant value. Additionally, clearance signals can be permanently set to TRUE when an LEA, as in the example of the PAL, is located at the end of the Logistics Line and does not need to wait for clearance from a downstream LEA.

\begin{graybox}
	\textbf{Best Practices for Choreography Configuration in Logistics Lines:}
	\begin{itemize}
		\item The Configurable Communication design pattern fundamentally enables the exchange of all information available in an LEA with other LEAs of a choreographed Logistics Line. However, to enable vendor-neutral interaction between LEAs, variables from the standardized MTP-based LEA interfaces according to Chapter~\ref{chap:Art1LEA} should be exchanged wherever possible.
		\item An exception is the interface of the External Lead. For this, \enquote{DoRequest and DoneReply relations} are recommended in accordance with \cite{Stutz.2026}, which enable the merging of different start and end points of a choreography relation across multiple LEA services.
	\end{itemize}
\end{graybox}

\subsection{Application Example: Transport PAL–Label–SH}
\label{subsec:LAAnwendung}

To illustrate the concepts described in Chapter~\ref{chap:Art3LA}, this section presents an exemplary application of the transport concept. Similarly to Chapter~\ref{chap:Art3LA}, a push transport from the PAL to SH1 is used as the example; however, a labeling LEA (LABEL) with a Processing node is additionally to be visited along the way. Figure~\ref{fig:TransportBeispiel} shows the LEAs and their transport nodes.

\begin{figure}[htbp]
	\centerline{\includegraphics[width=0.6\textwidth]{Inhalt/Abbildungen/08_Bereichs-Automatisierung/Transport_Anwendung/Beispiel.png}}
	\caption{Example of a push transport from a palletizer via a labeling unit to a stretch hood machine}
	\label{fig:TransportBeispiel}
\end{figure}

\subsubsection{0) LEA Integration}

As the basis for carrying out the transport process, the MTPs of the three LEAs are imported into the transport management in accordance with the mechanisms described in Section~\ref{subsec:Transportanbindung}, \emph{Transport Proxies} are generated, and these are connected to the LEA transport nodes. Figure~\ref{fig:IntegrationLea} shows the process of integrating the three LEAs; Figure~\ref{fig:IntegrationNode} shows the process of integrating all transport nodes available in the LEAs.\footnote{The individual TN proxies in the transport management are not shown here for clarity.}

\begin{figure}[htbp]
	\centerline{\includegraphics[width=0.8\textwidth]{Inhalt/Abbildungen/08_Bereichs-Automatisierung/Transport_Anwendung/00_Integration_LEA.png}}
	\caption{Integration of LEAs into a transport management system}
	\label{fig:IntegrationLea}
\end{figure}

\begin{figure}[htbp]
	\centerline{\includegraphics[width=1.0\textwidth]{Inhalt/Abbildungen/08_Bereichs-Automatisierung/Transport_Anwendung/00_Integration_Node.png}}
	\caption{Integration of all transport nodes of LEAs into a transport management system}
	\label{fig:IntegrationNode}
\end{figure}

\subsubsection{1) Transport Initiation}

Figure~\ref{fig:TransportInitiierung} shows the transport initiation process.

\begin{figure}[htbp]
	\centerline{\includegraphics[width=1.0\textwidth]{Inhalt/Abbildungen/08_Bereichs-Automatisierung/Transport_Anwendung/01_Transport_Initiierung.png}}
	\caption{Initiation of a push transport order}
	\label{fig:TransportInitiierung}
\end{figure}

In the context of the selected example, only the PAL has an Order node for initiating transport orders. Since no transport service is initially assigned to it, a new transport service (ID: transA) is created in the transport management and connected to the Order node (ID: 1) in accordance with the mechanisms described in Section~\ref{subsec:Transportanbindung}.

In the example, the PAL has a push transport requirement and starts the transport service in the \emph{PushRequested} procedure. The \emph{NextNode} is set to the InOutbound node of the PAL (ID: 2). The PAL obtains this ID by reading the \emph{ProxyId} from the \emph{TransportNodeManager} interface of its InOutbound node and sets it at the procedure parameter \emph{NextNode} of the transport service.

After the transport service is started, its OPC UA connection to the Order node is disconnected; however, it continues to run in the transport management. A new transport service (ID: transB) is created for the now-free Order node and assigned to it.

\subsubsection{2) Travel to the Start Node}

Figure~\ref{fig:FahrtStart} shows the process of traveling to the start node.

\begin{figure}[htbp]
	\centerline{\includegraphics[width=0.85\textwidth]{Inhalt/Abbildungen/08_Bereichs-Automatisierung/Transport_Anwendung/02_Transport_Fahrt_Start.png}}
	\caption{Travel to the start node}
	\label{fig:FahrtStart}
\end{figure}

The running transport service (ID: transA) is made available to a suitable FMS via its adapter. Depending on the FMS interface, this can occur at the initiative of the transport management or the FMS. Once the FMS has available capacity, it internally selects a suitable AGV to execute the transport process and reports the ID (here: agv3) back to an adapter. The adapter stores this ID in the \emph{ResourceId} variable of the transport service. Once the FMS begins executing the transport process, it restarts the transport service in the \emph{TransportRequested} procedure via its adapter. The selected AGV then travels to the InOutbound node of the PAL as the start node of the transport process.

\subsubsection{3) Coupling the Transport Service to the Start Node}

Figure~\ref{fig:AnbindungTransportdienst} shows the process of coupling the transport service to the start node.

\begin{figure}[htbp]
	\centerline{\includegraphics[width=0.96\textwidth]{Inhalt/Abbildungen/08_Bereichs-Automatisierung/Transport_Anwendung/03_Anbindung_Transportdienst.png}}
	\caption{Coupling the transport service to the start node}
	\label{fig:AnbindungTransportdienst}
\end{figure}

Shortly before the AGV arrives at the InOutbound node of the PAL, the FMS notifies its adapter of the imminent arrival. The adapter subsequently initiates a procedure change to \emph{TransportAwaitRequested} at the transport service. If the InOutbound node is not occupied by another transport service, the transport management assigns the transport service (ID: transA) to the \emph{Transport Proxy} of this node in accordance with Section~\ref{subsec:Transportanbindung}. The PAL then checks, based on the data from the transport service interface (e.g., \emph{ProductId}), whether it can process the transported product. For example, it may compare the \emph{ProductId} of the transport service with the \emph{ProductId} of the packaging order it is currently processing.

Depending on the result of this check, the PAL sets the procedure of the transport service to either \emph{TransportAwaited} or \emph{TransportDeclined}. In the event of a negative result (\emph{TransportDeclined}), the next transport node to be approached (procedure parameter \emph{NextNode}) must be redefined. To do so, the LEA may request an alternative node from the LOL or the operator via \emph{TransportNodeRequest} and reroute the transport accordingly. Alternatively, the FMS can direct the AGV to a safe position and the transport management places the transport service in the HELD state, signaling the need for manual intervention.

In the present example, a positive result is assumed. In this case, the PAL sets the procedure of the transport service to \emph{TransportAwaited}, which corresponds to the authorization for the AGV to approach the InOutbound node. When the FMS signals that the AGV has arrived at the node, the FMS adapter initiates a procedure change to \emph{TransportArrived}.

\subsubsection{4) Loading the AGV}

Figure~\ref{fig:BeladungAgv} shows the AGV loading process.

\begin{figure}[htbp]
	\centerline{\includegraphics[width=0.95\textwidth]{Inhalt/Abbildungen/08_Bereichs-Automatisierung/Transport_Anwendung/04_Beladung_AGV.png}}
	\caption{Loading the AGV}
	\label{fig:BeladungAgv}
\end{figure}

As soon as the InOutbound node of the PAL detects the \emph{TransportArrived} status, it initiates a procedure change to \emph{TransferFromLea} and the handover of the LO (here: pallet) to the AGV begins. Once the transfer is complete, the next transport node to be approached (here: the Processing node of the LABEL LEA) is determined in accordance with the mechanisms described in Section~\ref{subsec:ErmittlungNextNode} and set at the procedure parameter \emph{NextNode} of the transport service. The LEA then restarts the transport service in the \emph{TransferFromLeaSucceeded} procedure. The transport management recognizes this and decouples the transport service from the InOutbound node of the PAL, which is then again available for other transport orders.

\subsubsection{5) Travel to the Next Node Including Rerouting}

Figure~\ref{fig:FahrtNextNode} shows the process of traveling to the next node, including the option of rerouting.

\begin{figure}[htbp]
	\centerline{\includegraphics[width=0.85\textwidth]{Inhalt/Abbildungen/08_Bereichs-Automatisierung/Transport_Anwendung/05_Transport_Fahrt_NextNode_Rerouting.png}}
	\caption{Travel to the next node including rerouting}
	\label{fig:FahrtNextNode}
\end{figure}

As soon as the FMS adapter detects the \emph{TransferFromLeaSucceeded} procedure, it forwards the instruction to travel to the next transport node (here: the Processing node of the LABEL LEA) to the FMS. The FMS reports when the transport is underway, and the adapter initiates a procedure change to \emph{Transport}.

The transport management monitors the connections to the \emph{TransportClientManager} interfaces of the LEAs. If a connection loss is detected, or if a fault state is detected in the \emph{LeaStateCur} variable of the interface, rerouting of transport services that are en route to the affected faulty LEA can be initiated. In this case, an alternative transport node is determined internally within the LOL, initiated by the transport management, and set at the procedure parameter \emph{NextNode} of the transport service. The transport management then restarts the transport service in the \emph{Rerouting} procedure. In this way, the FMS adapter detects the rerouting and forwards the destination change to the FMS. Once the FMS has accepted the new destination, its adapter initiates a procedure change to \emph{Transport}.

\subsubsection{6) Processing}

Figure~\ref{fig:Processing} shows the process of handling an LO at a Processing node.

\begin{figure}[htbp]
	\centerline{\includegraphics[width=1.0\textwidth]{Inhalt/Abbildungen/08_Bereichs-Automatisierung/Transport_Anwendung/06_Transport_Processing.png}}
	\caption{Processing of a logistics object}
	\label{fig:Processing}
\end{figure}

If the approached transport node is a Processing node, the transport service is coupled to it following the same principle as described in Step 3. Processing of the LO then begins, which is signaled by a procedure change to \emph{Processing} initiated by the transport node of the LEA. In this case, \underline{no} handover of the LO to the LEA takes place; instead, it remains on the AGV. Once processing is complete, the next transport node to be approached is determined and set at the procedure parameter \emph{NextNode} of the transport service. The LEA then performs a procedure change to \emph{ProcessingSucceeded}. The transport management recognizes this and decouples the transport service from the Processing node. Subsequently, any number of additional Processing nodes can be approached in the same manner. Travel to the next transport node follows the same principle as described in Step 5.

\subsubsection{7) Unloading the AGV}

Figure~\ref{fig:EntladungAgv} shows the AGV unloading process.

\begin{figure}[htbp]
	\centerline{\includegraphics[width=1.0\textwidth]{Inhalt/Abbildungen/08_Bereichs-Automatisierung/Transport_Anwendung/07_Entladung_AGV.png}}
	\caption{Unloading the AGV}
	\label{fig:EntladungAgv}
\end{figure}

If no (further) Processing node is to be approached, the transport to the last transport node to be visited (destination node, here: the Inbound node of SH1) is coordinated following the same principle as described in Step 5. Shortly before the AGV arrives, the transport service is coupled to the transport node following the same principle as described in Step 3. The LO is then transferred from the AGV to the SH1, which is signaled by a procedure change to \emph{TransferToLea} initiated by the transport node of the LEA. After the successful transfer, the LEA triggers a procedure change to \emph{TransferToLeaSucceeded}. The transport management finally decouples the transport service from the transport node and terminates it via a \emph{Complete} command.
